\documentclass[a4paper,12pt]{article}

% ===========================================================================
%  Лабораторная работа 3.3.4 «Эффект Холла в полупроводниках»
%
%  Сборка:  cd 3.3.4/report && pdflatex main.tex && pdflatex main.tex
%  Таблицы читаются из ../data/processed/*.csv, итоговые числа --
%  из ../data/processed/results.tex, графики -- из images/*.pdf.
%  Всё это готовят ../scripts/process_data.py и ../scripts/plot_data.py.
%  До проведения измерений на месте чисел стоит «???», на месте
%  графиков -- рамки-заглушки.
% ===========================================================================

% ============================================================
% Общий преамбул для отчётов по лабораторным работам (labs_lpr)
% Подключается из report/main.tex командой:
%   % ============================================================
% Общий преамбул для отчётов по лабораторным работам (labs_lpr)
% Подключается из report/main.tex командой:
%   % ============================================================
% Общий преамбул для отчётов по лабораторным работам (labs_lpr)
% Подключается из report/main.tex командой:
%   \input{../../common/latex_template/preamble.tex}
% ============================================================

% Поддержка русского языка и кодировки
\usepackage[T2A]{fontenc}
\usepackage{lmodern}
\usepackage[utf8]{inputenc}
\usepackage[russian]{babel}

% Математические пакеты
\usepackage{amsmath, amssymb}

% Настройка полей
\usepackage{geometry}
\geometry{top=2cm, bottom=2cm, left=2.5cm, right=1.5cm}

% Колонтитулы
\usepackage{fancyhdr}
\usepackage{graphicx}

% Современные таблицы
\usepackage{tabularray}

% Работа с csv файлами
\usepackage{csvsimple-l3}
\usepackage{tikz}

% Графики
\usepackage{pgfplots}
\pgfplotsset{compat=1.18}

% Ссылки
% \usepackage{hyperref} % гипперссылки, нужно оформление вместо красного квадрата
\usepackage[russian]{cleveref} % умные ссылки, нужно переоформление для формул, рисунков и тд (вместо ф-л что-то нужно придумать)

% ============================================================

% Поддержка русского языка и кодировки
\usepackage[T2A]{fontenc}
\usepackage{lmodern}
\usepackage[utf8]{inputenc}
\usepackage[russian]{babel}

% Математические пакеты
\usepackage{amsmath, amssymb}

% Настройка полей
\usepackage{geometry}
\geometry{top=2cm, bottom=2cm, left=2.5cm, right=1.5cm}

% Колонтитулы
\usepackage{fancyhdr}
\usepackage{graphicx}

% Современные таблицы
\usepackage{tabularray}

% Работа с csv файлами
\usepackage{csvsimple-l3}
\usepackage{tikz}

% Графики
\usepackage{pgfplots}
\pgfplotsset{compat=1.18}

% Ссылки
% \usepackage{hyperref} % гипперссылки, нужно оформление вместо красного квадрата
\usepackage[russian]{cleveref} % умные ссылки, нужно переоформление для формул, рисунков и тд (вместо ф-л что-то нужно придумать)

% ============================================================

% Поддержка русского языка и кодировки
\usepackage[T2A]{fontenc}
\usepackage{lmodern}
\usepackage[utf8]{inputenc}
\usepackage[russian]{babel}

% Математические пакеты
\usepackage{amsmath, amssymb}

% Настройка полей
\usepackage{geometry}
\geometry{top=2cm, bottom=2cm, left=2.5cm, right=1.5cm}

% Колонтитулы
\usepackage{fancyhdr}
\usepackage{graphicx}

% Современные таблицы
\usepackage{tabularray}

% Работа с csv файлами
\usepackage{csvsimple-l3}
\usepackage{tikz}

% Графики
\usepackage{pgfplots}
\pgfplotsset{compat=1.18}

% Ссылки
% \usepackage{hyperref} % гипперссылки, нужно оформление вместо красного квадрата
\usepackage[russian]{cleveref} % умные ссылки, нужно переоформление для формул, рисунков и тд (вместо ф-л что-то нужно придумать)


% Пакеты, нужные именно этой работе
\usepackage{subcaption}
\usepackage{circuitikz}
\usepackage{siunitx}
\sisetup{output-decimal-marker={,}, group-digits=false}

% --- Служебные макросы ------------------------------------------------------
% Короткий пропуск (число/слово, которое надо вписать после измерений)
\newcommand{\qq}{\textbf{???}}
% Пропуск, требующий осмысленного текста
\newcommand{\fillme}{\textbf{?? НЕОБХОДИМО ЗАПОЛНИТЬ ??}}
% Имя файла моноширинным шрифтом (символы «_» и т.п. экранировать не нужно)
\newcommand{\filename}[1]{\texttt{\detokenize{#1}}}
% Число из CSV: десятичная запятая; пустая ячейка -> прочерк
% (protected: в «собираемых» таблицах csvsimple значения подставляются
%  через \csvexpval, а сам \cnum раскрывается уже внутри ячейки)
\NewDocumentCommand{\cnum}{m}{\edef\cnumtmp{#1}%
    \ifx\cnumtmp\empty --\else\num{\cnumtmp}\fi}
% График из report/images/. Если файла ещё нет -- рамка-заглушка.
% #1 -- имя файла, #2 -- доля ширины окружения.
\newcommand{\labimage}[2]{%
    \IfFileExists{images/#1}{%
        \includegraphics[width=#2\linewidth]{images/#1}%
    }{%
        \fbox{\parbox[c][0.55\dimexpr#2\linewidth\relax][c]{%
                \dimexpr#2\linewidth-2\fboxsep-2\fboxrule\relax}{%
            \centering\small\filename{images/#1}\\[2mm]
            (запустить \filename{scripts/plot_data.py})}}%
    }%
}

% --- Итоговые числа ----------------------------------------------------------
% Файл пишет process_data.py; если его нет, все величины -- «???».
\InputIfFileExists{../data/processed/results.tex}{}{}
\providecommand*{\ResK}{\qq}       \providecommand*{\ResKb}{\qq}
\providecommand*{\ResRH}{\qq}      \providecommand*{\ResRHabs}{\qq}
\providecommand*{\ResN}{\qq}       \providecommand*{\ResNcm}{\qq}
\providecommand*{\ResRtf}{\qq}     \providecommand*{\ResRho}{\qq}
\providecommand*{\ResRhoOhmCm}{\qq}\providecommand*{\ResSigma}{\qq}
\providecommand*{\ResMu}{\qq}      \providecommand*{\ResKrev}{\qq}
\providecommand*{\ResUohm}{\qq}    \providecommand*{\ResNu}{\qq}
\providecommand*{\ResVT}{\qq}      \providecommand*{\ResLambda}{\qq}
\providecommand*{\ResD}{\qq}       \providecommand*{\ResBmax}{\qq}
\providecommand*{\ResCalibN}{\qq}
\providecommand*{\ParA}{\qq}       \providecommand*{\ParH}{\qq}
\providecommand*{\ParL}{\qq}       \providecommand*{\ParSN}{\qq}
\providecommand*{\ParT}{\qq}       \providecommand*{\ParMeff}{\qq}
\providecommand*{\ParEpsB}{\qq}

% Настройка колонтитулов (как в остальных работах репозитория)
\pagestyle{fancy}
\fancyhf{}
\fancyhead[L]{%
    \IfFileExists{../../common/latex_template/miptlogo.png}{%
        \includegraphics[height=1cm]{../../common/latex_template/miptlogo.png}%
    }{%
        \textbf{[ЛОГОТИП МФТИ]}%
    }%
}
\fancyhead[R]{\textbf{Емельянов ДА, Байдин ФМ}}
\fancyfoot[C]{\thepage}
\renewcommand{\headrulewidth}{0.4pt}
\setlength{\headheight}{33pt}

% Подписи длинных таблиц (longtblr) по-русски
\DefTblrTemplate{contfoot-text}{default}{Продолжение на следующей странице}
\DefTblrTemplate{conthead-text}{default}{(продолжение)}

\begin{document}

% --- ТИТУЛЬНЫЙ ЛИСТ ---
\begin{titlepage}
\begin{center}
    {\large МОСКОВСКИЙ ФИЗИКО-ТЕХНИЧЕСКИЙ ИНСТИТУТ}
\end{center}
\begin{center}
    {\large Физтех-школа физики и исследований им. Ландау\break ЛФИ}
\end{center}

    \vspace{3cm}

\vspace{0.15cm}
{
\begin{center}
    \huge \bf Лабораторная работа №3.3.4 по курсу Общая физика\\
    \vspace{0.4cm}
    \LARGE Эффект Холла в полупроводниках.
\end{center}
}
\vspace{1.5cm}
\begin{flushright}
Емельянов Дмитрий Алексеевич\\
Байдин Фаддей Михайлович\\
\vspace{0.4cm}
Б02-502
\end{flushright}
\vspace{10cm}
\begin{center}
    \large \qq\ 2026, г. Долгопрудный
\end{center}
\end{titlepage}
% ----------------------

\section{Аннотация}
В работе исследуется эффект Холла в образце легированного германия в
геометрии мостика Холла. Через плоскую полупроводниковую пластинку,
помещённую в перпендикулярное ей магнитное поле электромагнита,
пропускается ток, и измеряется разность потенциалов между краями пластинки
в направлении, поперечном току. По зависимости ЭДС Холла от индукции поля
при различных токах через образец определяется постоянная Холла $R_H$, по её
знаку --- тип проводимости, а по величине --- концентрация $n$ основных
носителей заряда. Измерение удельного сопротивления образца в отсутствие
поля позволяет найти удельную проводимость $\sigma_0$ и подвижность $\mu$
носителей.

\section{Теоретические сведения}

\subsection{Проводимость и подвижность носителей}
Под действием внешнего электрического поля $\vec E$ носители заряда в
проводящей среде, многократно рассеиваясь на дефектах и колебаниях решётки,
приобретают среднюю установившуюся скорость дрейфа, пропорциональную полю:
\begin{equation}
    \label{eq:drift}
    \vec u_{\text{др}} = \mu\vec E,
\end{equation}
где коэффициент $\mu$ называется \emph{подвижностью} носителя тока. Знак
$\mu$ совпадает со знаком заряда $q$. Взаимодействие с решёткой можно
описать эффективной силой трения $\vec F_{\text{тр}} = -q\vec u/\mu$.
При концентрации носителей $n$ плотность тока $\vec j = qn\vec u = qn\mu\vec
E$, откуда следует закон Ома в дифференциальной форме $\vec j = \sigma\vec E$
с удельной проводимостью
\begin{equation}
    \label{eq:sigma}
    \sigma = qn\mu.
\end{equation}
В модели Друде, если $\nu$ --- частота столкновений носителя массы $m$ с
решёткой, подвижность равна
\begin{equation}
    \label{eq:drude}
    \mu = \frac{q}{m\nu}.
\end{equation}

В полупроводниках ток переносится электронами зоны проводимости и дырками
валентной зоны (дырке приписывают заряд $+e$). В легированном
полупроводнике, как правило, доминирует один тип носителей: электроны
($n$-тип) или дырки ($p$-тип).

\subsection{Эффект Холла}
Во внешнем магнитном поле $\vec B$ на движущийся носитель действует сила
Лоренца $\vec F = q\vec E + q\,\vec u\times\vec B$. Пусть ток течёт по оси $x$
вдоль плоской пластинки, а поле $\vec B$ направлено по оси $z$
перпендикулярно её плоскости. Магнитная составляющая силы Лоренца
отклоняет носители к одному из краёв пластинки; там накапливается заряд, и
возникает поперечное (\emph{холловское}) электрическое поле $E_y$. В
стационарном состоянии оно полностью компенсирует магнитную силу, так что
ток течёт строго вдоль $x$:
\begin{equation}
    \label{eq:Ey}
    E_y = u_x B = \frac{j_x}{nq}B.
\end{equation}

Уравнение баланса сил $q(\vec E + \vec u\times\vec B) - q\vec u/\mu = 0$
приводит к обобщённому закону Ома
\begin{equation}
    \label{eq:ohm_B}
    \vec E = \frac{\vec j}{\sigma_0} - \frac{1}{nq}\,\vec j\times\vec B,
    \qquad
    \hat\rho = \frac{1}{\sigma_0}
    \begin{pmatrix}
        1 & -\mu B & 0\\
        \mu B & 1 & 0\\
        0 & 0 & 1
    \end{pmatrix},
\end{equation}
где $\sigma_0 = qn\mu$ --- проводимость в отсутствие поля. Второе слагаемое
описывает эффект Холла. Отметим, что диагональная компонента
$\rho_{xx} = 1/\sigma_0$ от поля не зависит, т.\,е. в геометрии мостика
Холла (при одном типе носителей) продольное сопротивление образца не
меняется в магнитном поле.

\subsection{Мостик Холла}
Пусть пластинка имеет ширину $a$ (вдоль $y$), толщину $h$ (вдоль $z$), и
через неё течёт ток $I$, так что $j_x = I/(ah)$. Холловское напряжение между
краями пластинки $U_\perp = E_y a$ с учётом \cref{eq:Ey} равно
\begin{equation}
    \label{eq:hall}
    U_\perp = \frac{B}{nqh}\,I = R_H\,\frac{B I}{h},
\end{equation}
где
\begin{equation}
    \label{eq:RH}
    R_H = \frac{1}{nq}
\end{equation}
--- \emph{постоянная Холла}. Её знак совпадает со знаком заряда основных
носителей, поэтому по знаку ЭДС Холла при известных направлениях тока и
поля определяется тип проводимости. Модуль $R_H$ даёт концентрацию
носителей $n = 1/(e|R_H|)$. Из \cref{eq:sigma,eq:RH} следует выражение для
подвижности через измеряемые величины:
\begin{equation}
    \label{eq:mu}
    \mu = \sigma_0 R_H.
\end{equation}

Согласно \cref{eq:hall}, при фиксированном токе $I$ напряжение $U_\perp$
линейно зависит от $B$ с угловым коэффициентом
\begin{equation}
    \label{eq:k}
    k = \frac{dU_\perp}{dB} = \frac{R_H}{h}\,I,
\end{equation}
который, в свою очередь, пропорционален току. Поэтому, построив зависимость
$k(I)$ и определив её наклон $K = dk/dI$, получаем
\begin{equation}
    \label{eq:RH_from_K}
    R_H = K h.
\end{equation}

Удельное сопротивление материала находится по падению напряжения $U_{35}$
на участке образца длиной $l$ между контактами 3 и 5 в отсутствие поля:
\begin{equation}
    \label{eq:rho}
    \rho_0 = \frac{1}{\sigma_0} = \frac{U_{35}\,a h}{I\,l}.
\end{equation}

\subsection{Исключение омической добавки}
Холловские контакты 3 и 4 из-за неточности подпайки, как правило, лежат не
на одной эквипотенциали, поэтому измеряемое напряжение
$U_{34} = U_\perp + U_0$ содержит омическую добавку $U_0$, пропорциональную
току через образец и не зависящую от поля. Её можно исключить двумя
способами:
\begin{enumerate}
    \item измерить $U_0$ при выключенном магните для каждого тока $I$ и
          отсчитывать ЭДС Холла от этого значения:
          \begin{equation}
              \label{eq:U0}
              U_\perp = U_{34} - U_0;
          \end{equation}
    \item измерить $U_{34}$ при двух противоположных направлениях поля:
          при обращении поля ЭДС Холла меняет знак, а $U_0$ --- нет, так
          что
          \begin{equation}
              \label{eq:reverse}
              U_\perp = \frac{U_{34}^{(+)} - U_{34}^{(-)}}{2},
              \qquad
              U_0 = \frac{U_{34}^{(+)} + U_{34}^{(-)}}{2}.
          \end{equation}
\end{enumerate}
Основным в работе является первый способ, второй используется для контроля
при максимальном токе через образец.

\subsection{Кинетические оценки}
По найденной подвижности можно оценить (контрольный вопрос 6) частоту
столкновений $\nu = e/(m\mu)$ (\cref{eq:drude}), длину свободного пробега
$\lambda = \bar v/\nu$, где для невырожденного газа носителей
$\bar v \approx \sqrt{3k_BT/m}$, и коэффициент диффузии по соотношению
Эйнштейна
\begin{equation}
    \label{eq:einstein}
    D = \frac{\mu k_B T}{e}.
\end{equation}

\section{Используемое оборудование}
\begin{itemize}
    \item электромагнит с регулируемым источником питания, амперметр А1
          (ток $I_\text{М}$ через обмотки), переключатель направления тока
          К1;
    \item милливеберметр с пробной катушкой или миллитесламетр на основе
          датчика Холла (градуировка электромагнита);
    \item образец легированного германия в держателе с контактами 1--5;
    \item источник питания образца ($\simeq1{,}5$~В), реостат $R_2$,
          ключ К2, миллиамперметр А2 (ток $I$ через образец);
    \item цифровой вольтметр V (напряжения $U_{34}$ и $U_{35}$).
\end{itemize}
Параметры приборов и образца: \fillme.

\section{Экспериментальная установка и методика измерений}
Электрическая схема цепи образца показана на \cref{fig:setup}, расположение
контактов на образце --- на \cref{fig:sample}. Образец помещается в зазор
электромагнита так, что поле $\vec B$ перпендикулярно плоскости пластинки.
Ток через образец течёт между контактами 1 и 2 вдоль длинной стороны
пластинки, регулируется реостатом $R_2$ и не должен превышать $1$~мА.
Холловское напряжение снимается с контактов 3 и 4, расположенных на
противоположных краях пластинки, продольное падение напряжения --- с
контактов 3 и 5, расположенных на одном краю на расстоянии $l$ друг от
друга.

\begin{figure}[!ht]
    \centering
    \begin{circuitikz}[scale=0.95, transform shape]
        % образец
        \draw[fill=gray!15] (3,0) rectangle (7,1.6);
        \node at (5,0.8) {образец};
        \fill (3,0.8) circle (1.5pt) node[left=1pt, yshift=6pt] {\small 1};
        \fill (7,0.8) circle (1.5pt) node[right=1pt, yshift=6pt] {\small 2};
        \fill (4.5,1.6) circle (1.5pt) node[above left] {\small 3};
        \fill (4.5,0)   circle (1.5pt) node[below left] {\small 4};
        \fill (6.0,1.6) circle (1.5pt) node[above right] {\small 5};
        % цепь тока образца
        \node[jump crossing] (X) at (1.5,-1.0) {};
        \draw (3,0.8) -- (1.5,0.8) -- (X.north);
        \draw (X.south) -- (1.5,-2.2);
        \draw (1.5,-2.2) to[battery1, l_=$1{,}5$~В] (3.8,-2.2)
              to[switch, l_=К2] (5.6,-2.2)
              to[vR, l_=$R_2$] (7.6,-2.2)
              to[rmeter, t=mA, l_=А2] (8.6,-2.2) -- (8.6,0.8) -- (7,0.8);
        % вольтметр между 3 и 4
        \draw (4.5,1.6) -- (4.5,2.4) -- (0,2.4)
              to[rmeter, t=V] (0,-1.0) -- (X.west);
        \draw (X.east) -- (4.5,-1.0) -- (4.5,0);
        % поле
        \node at (8.0,1.9) {$\vec B\odot$};
    \end{circuitikz}
    \caption{Схема цепи образца. Вольтметр V подключается к контактам 3--4
             (ЭДС Холла) или 3--5 (продольное напряжение)}
    \label{fig:setup}
\end{figure}

\begin{figure}[!ht]
    \centering
    \begin{tikzpicture}[scale=1.1, >=stealth]
        \draw[fill=gray!15] (0,0) rectangle (6,1.8);
        \foreach \x/\y/\n/\p in {0/0.9/1/left, 6/0.9/2/right,
                                 2/1.8/3/above, 2/0/4/below, 4.5/1.8/5/above}
            {\fill (\x,\y) circle (1.6pt) node[\p] {\small\n};}
        \draw[->, thick] (0.5,0.9) -- (1.6,0.9) node[above] {$\vec j$};
        \draw[<->] (2,2.35) -- node[above] {$l$} (4.5,2.35);
        \draw[dashed] (2,1.85) -- (2,2.45) (4.5,1.85) -- (4.5,2.45);
        \draw[<->] (6.5,0) -- node[right] {$a$} (6.5,1.8);
        \node at (5.3,0.6) {$\vec B\odot$};
        \node[align=left, anchor=west] at (7.2,0.9) {толщина $h$\\(вдоль $\vec B$)};
    \end{tikzpicture}
    \caption{Образец и расположение контактов (вид вдоль поля)}
    \label{fig:sample}
\end{figure}

Порядок измерений:
\begin{enumerate}
    \item Градуировка электромагнита: зависимость индукции $B$ в зазоре от
          тока $I_\text{М}$ через обмотки (не менее 15 точек во всём
          диапазоне токов). При использовании милливеберметра измеряется
          поток $\Phi = BSN$ через пробную катушку с известным значением
          $SN$.
    \item Для каждого из 6--8 значений тока $I$ через образец:
          при выключенном магните измеряется $U_0$, затем снимается
          зависимость $U_{34}(I_\text{М})$ (6--8 точек). ЭДС Холла
          вычисляется по \cref{eq:U0}, индукция --- по градуировочной
          кривой (кусочно-линейная интерполяция).
    \item При максимальном токе через образец $U_{34}(I_\text{М})$
          измеряется также при обратном направлении поля (контроль по
          \cref{eq:reverse}).
    \item По направлениям тока, поля и знаку ЭДС Холла определяется знак
          основных носителей.
    \item Образец вынимается из зазора, при максимальном токе измеряется
          напряжение $U_{35}$; по \cref{eq:rho} вычисляется $\rho_0$.
\end{enumerate}
Угловые коэффициенты прямых $U_\perp(B)$ и $k(I)$ и их случайные
погрешности находятся методом наименьших квадратов (скрипт
\filename{scripts/process_data.py}). Погрешность $R_H$ включает также
погрешность толщины образца и систематическую погрешность градуировки поля
$\varepsilon_B$:
\begin{equation}
    \label{eq:err_RH}
    \frac{\sigma_{R_H}}{R_H} = \sqrt{\left(\frac{\sigma_K}{K}\right)^2
        + \left(\frac{\sigma_h}{h}\right)^2 + \varepsilon_B^2},
    \qquad
    \frac{\sigma_\mu}{\mu} = \sqrt{\left(\frac{\sigma_{R_H}}{R_H}\right)^2
        + \left(\frac{\sigma_{\rho}}{\rho_0}\right)^2}.
\end{equation}

\section{Результаты измерений и обработка данных}

\subsection{Параметры установки}
\begin{itemize}
    \item размеры образца: $a = \ParA$~мм, $h = \ParH$~мм,
          $l = \ParL$~мм;
    \item пробная катушка милливеберметра: $SN = \ParSN$~см$^2$
          (если градуировка проводилась тесламетром --- не используется);
    \item погрешность градуировки поля $\varepsilon_B = \ParEpsB$~\%;
    \item температура в лаборатории $T = \ParT$~К.
\end{itemize}
Все величины записаны в \filename{data/raw/setup_params.csv}.

\subsection{Градуировка электромагнита}
Результаты градуировки ($\ResCalibN$ точек) приведены в
\cref{tab:calibration} и на \cref{fig:calibration}. Максимальная индукция
в зазоре $B_{\max} = \ResBmax$~мТл.

\begin{table}[!ht]
    \centering
    \caption{Градуировка электромагнита}
    \label{tab:calibration}
    \csvreader[
        tabularray = {
            colspec = {|c|c|},
            hlines, vlines,
            row{1} = {font=\bfseries},
        },
        table head = {$I_\text{М}$, А & $B$, мТл \\},
    ]%
    {../data/processed/calibration_processed.csv}%
    {I_M_A=\cIM, B_mT=\cB}%
    {\cnum{\csvexpval\cIM} & \cnum{\csvexpval\cB}}
\end{table}

\begin{figure}[!ht]
    \centering
    \labimage{calibration.pdf}{0.75}
    \caption{Градуировочная кривая электромагнита $B(I_\text{М})$}
    \label{fig:calibration}
\end{figure}

Зависимость $B(I_\text{М})$ \fillme{} (линейна при малых токах,
наблюдается ли насыщение сердечника при больших).

\subsection{ЭДС Холла}
Измеренные значения $U_{34}$ и вычисленные по \cref{eq:U0} значения ЭДС
Холла приведены в \cref{tab:hall}; начальные напряжения $U_0$ для каждого
тока --- в \cref{tab:slopes}.

\csvreader[
    generic collected table = longtblr,
    generic table options = {[
            caption = {Зависимость ЭДС Холла от индукции поля при различных
                       токах через образец},
            label = {tab:hall},
        ]{
            colspec = {|c|c|c|c|c|c|},
            hlines, vlines,
            row{1} = {font=\bfseries},
        }},
    late after line = \\,
    table head = {№ & $I$, мА & $I_\text{М}$, А & $B$, мТл & $U_{34}$, мВ
                  & $U_\perp$, мВ \\},
]%
{../data/processed/hall_processed.csv}%
{series=\cS, I_mA=\cI, I_M_A=\cIM, B_mT=\cB, U34_mV=\cU, U_perp_mV=\cUh}%
{\csvexpval\cS & \cnum{\csvexpval\cI} & \cnum{\csvexpval\cIM} & \cnum{\csvexpval\cB} & \cnum{\csvexpval\cU} & \cnum{\csvexpval\cUh}}

Семейство зависимостей $U_\perp(B)$ показано на \cref{fig:family}. Все
зависимости \fillme{} (линейны / проходят через начало координат в пределах
погрешности). Угловые коэффициенты $k = dU_\perp/dB$, найденные МНК,
приведены в \cref{tab:slopes}.

\begin{figure}[!ht]
    \centering
    \labimage{hall_family.pdf}{0.85}
    \caption{Зависимость ЭДС Холла $U_\perp$ от индукции $B$ при различных
             токах $I$ через образец; прямые --- аппроксимация МНК}
    \label{fig:family}
\end{figure}

\begin{table}[!ht]
    \centering
    \caption{Начальные напряжения и угловые коэффициенты $k = dU_\perp/dB$
             зависимостей $U_\perp(B)$; $b$ --- свободный член прямой}
    \label{tab:slopes}
    \csvreader[
        tabularray = {
            colspec = {|c|c|c|c|c|c|c|},
            hlines, vlines,
            row{1} = {font=\bfseries},
        },
        table head = {№ & $I$, мА & $U_0$, мВ & $k$, мВ/Тл & $\sigma_k$, мВ/Тл
                      & $b$, мВ & $\sigma_b$, мВ \\},
    ]%
    {../data/processed/slopes.csv}%
    {series=\cS, I_mA=\cI, U0_mV=\cU, k_mV_T=\cK, sigma_k=\cSK,
     b_mV=\cBb, sigma_b=\cSB}%
    {\csvexpval\cS & \cnum{\csvexpval\cI} & \cnum{\csvexpval\cU} & \cnum{\csvexpval\cK} & \cnum{\csvexpval\cSK}
     & \cnum{\csvexpval\cBb} & \cnum{\csvexpval\cSB}}
\end{table}

\subsection{Постоянная Холла и концентрация носителей}
Зависимость углового коэффициента $k$ от тока через образец показана на
\cref{fig:k_vs_I}. Аппроксимация прямой $k = KI + k_0$ даёт
\begin{equation}
    K = \ResK~\frac{\text{В}}{\text{Тл}\cdot\text{А}},
    \qquad k_0 = \ResKb~\text{мВ/Тл}.
\end{equation}
Свободный член \fillme{} (в пределах погрешности равен нулю, как и следует
из \cref{eq:k}).

\begin{figure}[!ht]
    \centering
    \labimage{k_vs_I.pdf}{0.75}
    \caption{Зависимость углового коэффициента $k = dU_\perp/dB$ от тока
             через образец}
    \label{fig:k_vs_I}
\end{figure}

По \cref{eq:RH_from_K,eq:err_RH} постоянная Холла
\begin{equation}
    |R_H| = K h = \ResRHabs~\text{м}^3/\text{Кл},
\end{equation}
и концентрация основных носителей
\begin{equation}
    n = \frac{1}{e|R_H|} = \ResN~\text{м}^{-3}
      = \ResNcm~\text{см}^{-3}.
\end{equation}

\subsection{Контроль методом обращения поля}
При максимальном токе через образец ЭДС Холла измерена также по
\cref{eq:reverse} (\cref{tab:reverse}, \cref{fig:reverse}). Угловой
коэффициент
\begin{equation}
    k_{\text{обр}} = \ResKrev~\text{мВ/Тл}
\end{equation}
\fillme{} (сравнить с $k$ для того же тока из \cref{tab:slopes}).
Полусумма $\bigl(U_{34}^{(+)} + U_{34}^{(-)}\bigr)/2 = \ResUohm$~мВ
\fillme{} (сравнить с $U_0$ при том же токе).

\begin{table}[!ht]
    \centering
    \caption{Измерение ЭДС Холла с обращением направления поля}
    \label{tab:reverse}
    \csvreader[
        tabularray = {
            colspec = {|c|c|c|c|c|c|},
            hlines, vlines,
            row{1} = {font=\bfseries},
        },
        table head = {$I_\text{М}$, А & $|B|$, мТл & $U_{34}^{(+)}$, мВ
                      & $U_{34}^{(-)}$, мВ & $U_\perp$, мВ & $U_0$, мВ \\},
    ]%
    {../data/processed/reverse_processed.csv}%
    {I_M_A=\cIM, B_mT=\cB, U34_plus_mV=\cUp, U34_minus_mV=\cUm,
     U_perp_mV=\cUh, U_ohm_mV=\cUo}%
    {\cnum{\csvexpval\cIM} & \cnum{\csvexpval\cB} & \cnum{\csvexpval\cUp} & \cnum{\csvexpval\cUm} & \cnum{\csvexpval\cUh}
     & \cnum{\csvexpval\cUo}}
\end{table}

\begin{figure}[!ht]
    \centering
    \labimage{reverse.pdf}{0.75}
    \caption{Напряжение $U_{34}$ при двух направлениях поля и ЭДС Холла,
             вычисленная по \cref{eq:reverse}}
    \label{fig:reverse}
\end{figure}

\subsection{Тип проводимости}
Ток через образец течёт от контакта \qq{} к контакту \qq; поле при
положении 1 разъёма К1 направлено \fillme; потенциал контакта 3 при этом
\qq{} потенциала контакта 4 ($U_\perp$ \qq~0). Сила Лоренца
$q\,\vec u\times\vec B$ отклоняет носители к краю \qq, следовательно,
основные носители заряда --- \fillme{} (электроны / дырки), и образец
обладает \fillme{} ($n$- / $p$-) типом проводимости; $R_H$ \qq~0.

% Рисунок из тетради: направления тока, поля и отклонения носителей.
% \begin{figure}[!ht]
%     \centering
%     \labimage{carrier_sign.png}{0.6}
%     \caption{Направления тока, поля и отклонение носителей в образце}
%     \label{fig:sign}
% \end{figure}

\subsection{Удельное сопротивление, проводимость и подвижность}
При удалённом из зазора образце измерено напряжение между контактами 3 и 5
(\cref{tab:resistivity}).

\begin{table}[!ht]
    \centering
    \caption{Измерение удельного сопротивления (без поля)}
    \label{tab:resistivity}
    \csvreader[
        tabularray = {
            colspec = {|c|c|c|},
            hlines, vlines,
            row{1} = {font=\bfseries},
        },
        table head = {$I$, мА & $U_{35}$, мВ & $R_{35} = U_{35}/I$, Ом \\},
    ]%
    {../data/processed/resistivity_processed.csv}%
    {I_mA=\cI, U35_mV=\cU, R35_Ohm=\cR}%
    {\cnum{\csvexpval\cI} & \cnum{\csvexpval\cU} & \cnum{\csvexpval\cR}}
\end{table}

Сопротивление участка 3--5 $R_{35} = \ResRtf$~Ом. По \cref{eq:rho}
\begin{equation}
    \rho_0 = R_{35}\frac{ah}{l} = \ResRho~\text{Ом}\cdot\text{м}
           = \ResRhoOhmCm~\text{Ом}\cdot\text{см},
    \qquad
    \sigma_0 = \frac{1}{\rho_0} = \ResSigma~\text{См/м}.
\end{equation}
Подвижность носителей по \cref{eq:mu}:
\begin{equation}
    \mu = \sigma_0|R_H| = \ResMu~\frac{\text{см}^2}{\text{В}\cdot\text{с}}.
\end{equation}

\subsection{Оценка кинетических характеристик носителей}
Принимая эффективную массу носителей $m = \ParMeff\,m_e$ и
$T = \ParT$~К, получаем:
\begin{itemize}
    \item частота столкновений $\nu = e/(m\mu) = \ResNu$~с$^{-1}$
          (время свободного пробега $\tau = 1/\nu$);
    \item средняя тепловая скорость $\bar v \approx \sqrt{3k_BT/m} = \ResVT$~м/с;
    \item длина свободного пробега $\lambda = \bar v/\nu = \ResLambda$~нм;
    \item коэффициент диффузии (\cref{eq:einstein})
          $D = \mu k_BT/e = \ResD$~см$^2$/с.
\end{itemize}
Оценки носят порядковый характер: эффективные массы носителей в германии
отличаются от $m_e$ (для электронов $m^*\approx0{,}12\,m_e$, для дырок
$m^*\approx0{,}3\,m_e$ --- в усреднённом по направлениям смысле),
соответственно $\nu$ возрастает, а $\lambda$ уменьшается в
$\sqrt{m_e/m^*}$ раз.

\section{Обсуждение результатов}
\begin{itemize}
    \item Зависимости $U_\perp(B)$ при всех токах \fillme{} (линейны), что
          подтверждает \cref{eq:hall}; наклоны $k$ пропорциональны току
          через образец (\cref{fig:k_vs_I}).
    \item Постоянная Холла $|R_H| = \ResRHabs$~м$^3$/Кл, концентрация
          носителей $n = \ResNcm$~см$^{-3}$. Для сравнения, собственная
          концентрация носителей в германии при комнатной температуре
          $n_i\approx2{,}4\cdot10^{13}$~см$^{-3}$; полученное значение
          \fillme{} (во сколько раз больше), что подтверждает примесный
          характер проводимости.
    \item Тип проводимости --- \fillme.
    \item Подвижность $\mu = \ResMu$~см$^2$/(В$\cdot$с). Табличные
          значения для чистого германия при $300$~К: для электронов
          $\approx3900$, для дырок $\approx1900$~см$^2$/(В$\cdot$с).
          \fillme{} (сравнить с табличным значением для найденного типа
          носителей; в легированном образце подвижность может быть ниже
          из-за рассеяния на ионизированных примесях).
    \item Метод обращения поля даёт $k_{\text{обр}} = \ResKrev$~мВ/Тл,
          \fillme{} (согласуется / не согласуется с основным методом).
    \item Основные источники погрешности: \fillme{} (градуировка поля
          $\varepsilon_B$, точность размеров образца, дрейф $U_0$ из-за
          нагрева образца током, неоднородность поля в зазоре).
\end{itemize}

\section{Заключение}
В работе исследован эффект Холла в образце легированного германия:
\begin{itemize}
    \item проведена градуировка электромагнита, максимальное поле в зазоре
          $B_{\max} = \ResBmax$~мТл;
    \item подтверждена линейная зависимость ЭДС Холла от индукции поля и
          тока через образец;
    \item определена постоянная Холла $|R_H| = \ResRHabs$~м$^3$/Кл и тип
          проводимости образца --- \fillme;
    \item найдена концентрация основных носителей
          $n = \ResN$~м$^{-3}$;
    \item измерены удельное сопротивление
          $\rho_0 = \ResRho$~Ом$\cdot$м и проводимость
          $\sigma_0 = \ResSigma$~См/м;
    \item вычислена подвижность носителей
          $\mu = \ResMu$~см$^2$/(В$\cdot$с) и оценены длина свободного
          пробега $\lambda \sim \ResLambda$~нм и коэффициент диффузии
          $D \sim \ResD$~см$^2$/с.
\end{itemize}

\end{document}
